\chapter{ARIMA and Unit Roots}
\label{cap:arima}
% ── index ───────────────────────────────────────────────────────
\index{ARIMA|textbf}
\index{unit root|textbf}
\index{Dickey-Fuller test|textbf}
\index{ADF|see{Dickey-Fuller test}}
\index{adf@\texttt{adf()}|textbf}
\index{order of integration}
\index{KPSS test}
\index{kpss@\texttt{kpss()}}

% ─────────────────────────────────────────────────────────────────────────────
\section{Integrated processes}

A process $\{y_t\}$ is \textbf{integrated of order $d$} --- denoted $I(d)$
--- when it needs to be differenced $d$ times to become stationary.
The most common case in economics is $I(1)$: the first difference is stationary,
but the level is not.

The simplest example of an $I(1)$ process is the \textbf{random walk}:
\[
  y_t = y_{t-1} + \varepsilon_t, \qquad \varepsilon_t \sim \mathrm{WN}(0, \sigma^2)
\]

The variance grows with $t$ ($\mathrm{Var}[y_t] = t\sigma^2$), so the
process is not stationary. The first difference $\Delta y_t = y_t -
y_{t-1} = \varepsilon_t$ is white noise --- therefore $I(0)$.

The random walk \emph{with drift} adds a constant $\mu$:
\[
  y_t = \mu + y_{t-1} + \varepsilon_t
\]
with deterministic trend $E[y_t] = y_0 + \mu t$ --- the basic model for
asset prices, exchange rates and many macroeconomic series.

% ─────────────────────────────────────────────────────────────────────────────
\section{Augmented Dickey-Fuller Test}

The ADF (\emph{Augmented Dickey-Fuller}) test formally tests for the presence of
a unit root. The auxiliary model estimated is:
\[
  \Delta y_t = \alpha + \delta y_{t-1}
               + \sum_{j=1}^{k} \gamma_j \Delta y_{t-j} + \varepsilon_t
\]

with hypotheses:
\[
  H_0: \delta = 0 \quad \text{(unit root)} \qquad
  H_1: \delta < 0 \quad \text{(stationary)}
\]

The $t$-statistic of $\hat\delta$ does not follow the standard $t$-distribution under
$H_0$ --- the Dickey-Fuller distribution with tabulated critical values is used.

\begin{aviso}
  ADF critical values are \textbf{negative}: $H_0$ is rejected
  when the statistic is sufficiently \emph{negative}. A positive value
  (or near zero) does not reject $H_0$ --- evidence of a unit root.
\end{aviso}

% ─────────────────────────────────────────────────────────────────────────────
\section{The ARIMA model}

ARIMA($p$,$d$,$q$) generalizes ARMA for nonstationary series,
applying $d$ differences before estimation:

\[
  \underbrace{\phi(L)(1-L)^d y_t}_{\text{AR in differences}}
  = c + \underbrace{\theta(L)\varepsilon_t}_{\text{MA}}
\]

where $L$ is the lag operator ($Ly_t = y_{t-1}$) and $(1-L)^d$ is the
difference operator of order $d$.

\begin{center}
\begin{tabular}{lll}
\toprule
\textbf{Model} & \textbf{Notation} & \textbf{Process} \\
\midrule
Random walk & ARIMA(0,1,0) & $\Delta y_t = c + \varepsilon_t$ \\
AR(1) in differences & ARIMA(1,1,0) & $\Delta y_t = c + \phi\,\Delta y_{t-1} + \varepsilon_t$ \\
MA(1) in differences & ARIMA(0,1,1) & $\Delta y_t = c + \varepsilon_t + \theta\,\varepsilon_{t-1}$ \\
\bottomrule
\end{tabular}
\end{center}

% ─────────────────────────────────────────────────────────────────────────────
\section{DGP Verification}

\subsection*{Random walk: ADF does not reject; AR(1): ADF rejects}

The DGP contrasts the two cases: random walk (I(1)) generated by
\lstinline{cumsum(e)} versus stationary AR(1) ($\phi=0{,}7$) generated by
a loop.

\begin{lstlisting}[caption={ADF on random walk vs.\ AR(1)}]
seed(42)
input df
id
1
end
for i in 1..=9 { df = append(df, df) }
df |> filter(_n <= 500)
generate df e  = rnormal()
generate df rw = cumsum(e)
generate df y  = 0.0
generate df t  = _n
let n = nrow(df)
let i = 2
while i <= n {
    let ly = mean(df, y, if = t == i - 1)
    let et = mean(df, e, if = t == i)
    replace df y = 0.7 * ly + et if t == i
    i = i + 1
}
adf(df, rw)
adf(df, y)
\end{lstlisting}

\begin{lstlisting}[language={}, basicstyle=\ttfamily\small,
  frame=none, numbers=none, backgroundcolor=\color{codbg},
  xleftmargin=0pt, xrightmargin=0pt, breaklines=false]
=============== Augmented Dickey-Fuller Test ===============
  Variable: rw   Lags used: 7
  H₀: series has a unit root (non-stationary)
  Test statistic:     -0.0980
  Critical values:  1%=-3.430  5%=-2.860  10%=-2.570
  Conclusion: Does not reject H₀ — unit root present

=============== Augmented Dickey-Fuller Test ===============
  Variable: y   Lags used: 7
  H₀: series has a unit root (non-stationary)
  Test statistic:     -6.4616
  Critical values:  1%=-3.430  5%=-2.860  10%=-2.570
  Conclusion: REJECTS H₀ — stationary
\end{lstlisting}

The random walk has statistic $-0{,}10$ --- above the negative critical values, does not
reject $H_0$. The AR(1) with $\phi=0{,}7$
has statistic $-6{,}46$ --- far below the 1\% critical value ($-3{,}43$),
rejects $H_0$ with high confidence.

\subsection*{ARIMA(0,1,0) estimation and order selection}

With the unit root identified, estimate with $d=1$. Three competing orders
for the true random walk:

\begin{lstlisting}[caption={Order selection for I(1) series}]
let m010 = arima(df, rw, p=0, d=1, q=0)
let m110 = arima(df, rw, p=1, d=1, q=0)
let m011 = arima(df, rw, p=0, d=1, q=1)
display m010
display m110
display m011
\end{lstlisting}

\begin{lstlisting}[language={}, basicstyle=\ttfamily\small,
  frame=none, numbers=none, backgroundcolor=\color{codbg},
  xleftmargin=0pt, xrightmargin=0pt, breaklines=false]
-- ARIMA(0,1,0): AIC=156.77  BIC=160.97
   intercept  0.4863  std=0.0128  z=38.14  p=0.000

-- ARIMA(1,1,0): AIC=158.70  BIC=167.10
   intercept  0.4805  std=0.0253  z=18.96  p=0.000
   ar.L1      0.0120  std=0.0450  z=0.266  p=0.791  [insig.]

-- ARIMA(0,1,1): AIC=157.80  BIC=166.19
   intercept  0.4855  std=0.0128  z=38.07  p=0.000
   ma.L1      0.0058  std=0.0453  z=0.128  p=0.898  [insig.]
\end{lstlisting}

The true model ARIMA(0,1,0) has the lowest AIC and BIC. The additional AR and MA
terms are insignificant ($p > 0{,}7$) --- BIC, which penalizes complexity more
strongly, confirms the parsimony of the simple model.

% ─────────────────────────────────────────────────────────────────────────────
\section{Workflow}

The standard protocol for univariate time series:

\begin{enumerate}
  \item \textbf{Inspect} the level of the series --- visible trend suggests
        nonstationarity.
  \item \textbf{Test} with ADF --- failing to reject $H_0$ indicates $d \geq 1$.
  \item \textbf{Difference} once and re-test with ADF on the differenced series
        --- confirms $d=1$.
  \item \textbf{Select} $(p,q)$ by AIC/BIC grid over the differenced series.
  \item \textbf{Validate} residuals: they should be white noise.
\end{enumerate}

% ─────────────────────────────────────────────────────────────────────────────
\section{Syntax}

\begin{lstlisting}
adf(df, y)
adf(df, y, lags=4)
arima(df, y, p=0, d=1, q=0)
arima(df, y, p=1, d=1, q=1)
\end{lstlisting}

\begin{lstlisting}[caption={Typical workflow with real data}]
load "exchange.csv" as df

adf(df, rate)

// if H0 is not rejected: estimate with d=1
let m = arima(df, rate, p=1, d=1, q=0)
display m
let yhat = predict(m, "xb")
\end{lstlisting}

\section{Empirical validation}

The random walk DGP in this chapter is used in the \texttt{arima\_book}
validation case in \texttt{validation/cases/}. The case estimates an
ARIMA(1,1,0) in \hayashi{} and compares it with reference implementations in
R and Python. Run \lstinline{hay validate} to see the current status.
